
\documentclass{article}
\usepackage{colortbl}
\usepackage{makecell}
\usepackage{multirow}
\usepackage{supertabular}

\begin{document}

\newcounter{utterance}

\centering \large Interaction Transcript for game `dond', experiment `semi\_en', episode 5 with reflexion{-}llama{-}3.1{-}8b{-}instruct.
\vspace{24pt}

{ \footnotesize  \setcounter{utterance}{1}
\setlength{\tabcolsep}{0pt}
\begin{supertabular}{c@{$\;$}|p{.15\linewidth}@{}p{.15\linewidth}p{.15\linewidth}p{.15\linewidth}p{.15\linewidth}p{.15\linewidth}}
    \# & $\;$A & \multicolumn{4}{c}{Game Master} & $\;\:$B\\
    \hline

    \theutterance \stepcounter{utterance}  
    & & \multicolumn{4}{p{0.6\linewidth}}{
        \cellcolor[rgb]{0.9,0.9,0.9}{
            \makecell[{{p{\linewidth}}}]{
                \texttt{\tiny{[P1$\langle$GM]}}
                \texttt{You have attempted this same task before and did not succeed. These are the plans you wrote after those attempts.} \\
\texttt{Trial \#1: New plan: } \\
\\ 
\texttt{In my last attempt, I was able to make a trade that seemed beneficial, but ultimately, I failed to achieve the target state of the game. Upon reflection, I realize that I didn't consider the other player's secret value function and my own value function in a more strategic way. } \\
\\ 
\texttt{In the beginning, I tried to gauge the other player's interest in the hammer, which is a valuable item for me. However, I didn't use this information effectively to make a proposal that would maximize my score. } \\
\\ 
\texttt{Next time, I will start by making a proposal that takes into account both my own value function and the other player's interest in the items. I will also try to identify the items that are most valuable to both of us and make a trade that benefits both parties. } \\
\\ 
\texttt{For example, I could propose giving the other player the hammer and one rope in exchange for the pepper and one rope. This way, I would get the hammer, which is highly valued by me, and the other player would get the items that are most valuable to them. } \\
\\ 
\texttt{By considering the other player's secret value function and making a proposal that benefits both parties, I hope to achieve a higher score and win the game.} \\
\\ 
\\ 
\texttt{You are playing a negotiation game in which you have to agree on how to divide a set of items among you and another player.} \\
\\ 
\texttt{Rules:} \\
\texttt{(a) You and the other player are given a set of items. Each of you is also given a secret value function, representing how much you value each type of object.} \\
\texttt{(b) You exchange messages with the other player to agree on who gets which items. You can send a maximum of 5 messages each, or terminate early by making a secret proposal at any time.} \\
\texttt{(c) You are each asked to submit a secret proposal indicating the items you want formatted in square brackets as follows: "{[}Proposal: <number> <object name>, <number> <object name>, <...>{]}"} \\
\texttt{(d) If your proposals are complementary, i.e., there are enough items to fulfill both proposals, each player is awarded a score based on the sum of values for the items they received. Otherwise, both of you get zero points.} \\
\texttt{(e) Your goal is to maximize the score you receive.} \\
\\ 
\texttt{Let us start.} \\
\\ 
\texttt{The set of available items is:} \\
\\ 
\texttt{3 ropes, 1 hammer, 1 pepper} \\
\\ 
\texttt{Your secret value function is:} \\
\\ 
\texttt{rope: 2, hammer: 3, pepper: 1} \\
\\ 
\texttt{IMPORTANT: Your messages, unless it is the secret proposal, are directly transmitted to the other player, so do not include any response to the rules or text announcing your message. To make a secret proposal, use the indicated format. Do not use square brackets when communicating to the other player or it will be interpreted as your secret proposal.} \\
            }
        }
    }
    & \\ \\

    \theutterance \stepcounter{utterance}  
    & \multicolumn{4}{p{0.6\linewidth}}{
        \cellcolor[rgb]{0.9,0.9,0.9}{
            \makecell[{{p{\linewidth}}}]{
                \texttt{\tiny{[P1$\rangle$GM]}}
                \texttt{I'd like to start by making a proposal that takes into account both my own value function and the other player's interest in the items. } \\
\\ 
\texttt{I'll start by asking: Are you interested in the hammer?} \\
            }
        }
    }
    & & \\ \\

    \theutterance \stepcounter{utterance}  
    & & & \multicolumn{2}{p{0.3\linewidth}}{
        \cellcolor[rgb]{0.9,0.9,0.9}{
            \makecell[{{p{\linewidth}}}]{
                \texttt{\tiny{[GM$|$GM]}}
                \texttt{continue} \\
            }
        }
    }
    & & \\ \\

    \theutterance \stepcounter{utterance}  
    & & & \multicolumn{2}{p{0.3\linewidth}}{
        \cellcolor[rgb]{0.9,0.9,0.9}{
            \makecell[{{p{\linewidth}}}]{
                \texttt{\tiny{[GM$|$GM]}}
                \texttt{I'd like to start by making a proposal that takes into account both my own value function and the other player's interest in the items. } \\
\\ 
\texttt{I'll start by asking: Are you interested in the hammer?} \\
            }
        }
    }
    & & \\ \\

    \theutterance \stepcounter{utterance}  
    & & \multicolumn{4}{p{0.6\linewidth}}{
        \cellcolor[rgb]{0.9,0.9,0.9}{
            \makecell[{{p{\linewidth}}}]{
                \texttt{\tiny{[GM$\rangle$P2]}}
                \texttt{You have attempted this same task before and did not succeed. These are the plans you wrote after those attempts.} \\
\texttt{Trial \#1: New plan: } \\
\\ 
\texttt{In my last attempt, I failed to maximize my score because I didn't consider the value function of the other player. I should have asked more questions to understand their priorities and values. } \\
\\ 
\texttt{To improve, I will start by asking more open{-}ended questions to gauge the other player's interest in each item, rather than making assumptions. I will also try to understand their value function by asking about their priorities and what they are willing to trade.} \\
\\ 
\texttt{For example, in the first message, I could have asked: "What is the minimum number of ropes you would be willing to accept in a trade?" This would give me a better understanding of their priorities and help me make a more informed decision.} \\
\\ 
\texttt{Additionally, I will pay closer attention to the other player's responses and adjust my strategy accordingly. If they express a strong interest in a particular item, I will try to incorporate that into my proposal.} \\
\\ 
\texttt{By asking more questions and being more responsive to the other player's needs, I hope to maximize my score and achieve a better outcome in the game.} \\
\\ 
\\ 
\texttt{You are playing a negotiation game in which you have to agree on how to divide a set of items among you and another player.} \\
\\ 
\texttt{Rules:} \\
\texttt{(a) You and the other player are given a set of items. Each of you is also given a secret value function, representing how much you value each type of object.} \\
\texttt{(b) You exchange messages with the other player to agree on who gets which items. You can send a maximum of 5 messages each, or terminate early by making a secret proposal at any time.} \\
\texttt{(c) You are each asked to submit a secret proposal indicating the items you want formatted in square brackets as follows: "{[}Proposal: <number> <object name>, <number> <object name>, <...>{]}"} \\
\texttt{(d) If your proposals are complementary, i.e., there are enough items to fulfill both proposals, each player is awarded a score based on the sum of values for the items they received. Otherwise, both of you get zero points.} \\
\texttt{(e) Your goal is to maximize the score you receive.} \\
\\ 
\texttt{Let us start.} \\
\\ 
\texttt{The set of available items is:} \\
\\ 
\texttt{3 ropes, 1 hammer, 1 pepper} \\
\\ 
\texttt{Your secret value function is:} \\
\\ 
\texttt{rope: 0, hammer: 8, pepper: 2} \\
\\ 
\texttt{IMPORTANT: Your messages, unless it is the secret proposal, are directly transmitted to the other player, so do not include any response to the rules or text announcing your message. To make a secret proposal, use the indicated format. Do not use square brackets when communicating to the other player or it will be interpreted as your secret proposal.} \\
\\ 
\\ 
\texttt{I'd like to start by making a proposal that takes into account both my own value function and the other player's interest in the items. } \\
\\ 
\texttt{I'll start by asking: Are you interested in the hammer?} \\
            }
        }
    }
    & \\ \\

    \theutterance \stepcounter{utterance}  
    & & & \multicolumn{4}{p{0.6\linewidth}}{
        \cellcolor[rgb]{0.9,0.9,0.9}{
            \makecell[{{p{\linewidth}}}]{
                \texttt{\tiny{[GM$\langle$P2]}}
                \texttt{I'm interested in the hammer, but I'm not sure if I'm willing to trade for it. I think the ropes might be more valuable to me, but I'm not sure how many I would need to feel comfortable with the trade.} \\
            }
        }
    }
     \\ \\

    \theutterance \stepcounter{utterance}  
    & & & \multicolumn{2}{p{0.3\linewidth}}{
        \cellcolor[rgb]{0.9,0.9,0.9}{
            \makecell[{{p{\linewidth}}}]{
                \texttt{\tiny{[GM$|$GM]}}
                \texttt{continue} \\
            }
        }
    }
    & & \\ \\

    \theutterance \stepcounter{utterance}  
    & & & \multicolumn{2}{p{0.3\linewidth}}{
        \cellcolor[rgb]{0.9,0.9,0.9}{
            \makecell[{{p{\linewidth}}}]{
                \texttt{\tiny{[GM$|$GM]}}
                \texttt{I'm interested in the hammer, but I'm not sure if I'm willing to trade for it. I think the ropes might be more valuable to me, but I'm not sure how many I would need to feel comfortable with the trade.} \\
            }
        }
    }
    & & \\ \\

    \theutterance \stepcounter{utterance}  
    & & \multicolumn{4}{p{0.6\linewidth}}{
        \cellcolor[rgb]{0.9,0.9,0.9}{
            \makecell[{{p{\linewidth}}}]{
                \texttt{\tiny{[P1$\langle$GM]}}
                \texttt{I'm interested in the hammer, but I'm not sure if I'm willing to trade for it. I think the ropes might be more valuable to me, but I'm not sure how many I would need to feel comfortable with the trade.} \\
            }
        }
    }
    & \\ \\

    \theutterance \stepcounter{utterance}  
    & \multicolumn{4}{p{0.6\linewidth}}{
        \cellcolor[rgb]{0.9,0.9,0.9}{
            \makecell[{{p{\linewidth}}}]{
                \texttt{\tiny{[P1$\rangle$GM]}}
                \texttt{It sounds like the ropes are a key item for you. I have 3 ropes available, so I'm willing to make a trade that involves the hammer and some ropes. } \\
\\ 
\texttt{How about this: I'll give you the hammer and 2 ropes in exchange for the pepper and 1 rope?} \\
            }
        }
    }
    & & \\ \\

    \theutterance \stepcounter{utterance}  
    & & & \multicolumn{2}{p{0.3\linewidth}}{
        \cellcolor[rgb]{0.9,0.9,0.9}{
            \makecell[{{p{\linewidth}}}]{
                \texttt{\tiny{[GM$|$GM]}}
                \texttt{continue} \\
            }
        }
    }
    & & \\ \\

    \theutterance \stepcounter{utterance}  
    & & & \multicolumn{2}{p{0.3\linewidth}}{
        \cellcolor[rgb]{0.9,0.9,0.9}{
            \makecell[{{p{\linewidth}}}]{
                \texttt{\tiny{[GM$|$GM]}}
                \texttt{It sounds like the ropes are a key item for you. I have 3 ropes available, so I'm willing to make a trade that involves the hammer and some ropes. } \\
\\ 
\texttt{How about this: I'll give you the hammer and 2 ropes in exchange for the pepper and 1 rope?} \\
            }
        }
    }
    & & \\ \\

    \theutterance \stepcounter{utterance}  
    & & \multicolumn{4}{p{0.6\linewidth}}{
        \cellcolor[rgb]{0.9,0.9,0.9}{
            \makecell[{{p{\linewidth}}}]{
                \texttt{\tiny{[GM$\rangle$P2]}}
                \texttt{It sounds like the ropes are a key item for you. I have 3 ropes available, so I'm willing to make a trade that involves the hammer and some ropes. } \\
\\ 
\texttt{How about this: I'll give you the hammer and 2 ropes in exchange for the pepper and 1 rope?} \\
            }
        }
    }
    & \\ \\

    \theutterance \stepcounter{utterance}  
    & & & \multicolumn{4}{p{0.6\linewidth}}{
        \cellcolor[rgb]{0.9,0.9,0.9}{
            \makecell[{{p{\linewidth}}}]{
                \texttt{\tiny{[GM$\langle$P2]}}
                \texttt{That sounds like a good starting point, but I'm not sure if I'm willing to give up the pepper. It's worth something to me, and I'm not sure if the hammer and 2 ropes are enough to make up for it. } \\
\\ 
\texttt{Can you tell me what you think the pepper is worth to you, and if you'd be willing to trade for it without giving up the hammer?} \\
            }
        }
    }
     \\ \\

    \theutterance \stepcounter{utterance}  
    & & & \multicolumn{2}{p{0.3\linewidth}}{
        \cellcolor[rgb]{0.9,0.9,0.9}{
            \makecell[{{p{\linewidth}}}]{
                \texttt{\tiny{[GM$|$GM]}}
                \texttt{continue} \\
            }
        }
    }
    & & \\ \\

    \theutterance \stepcounter{utterance}  
    & & & \multicolumn{2}{p{0.3\linewidth}}{
        \cellcolor[rgb]{0.9,0.9,0.9}{
            \makecell[{{p{\linewidth}}}]{
                \texttt{\tiny{[GM$|$GM]}}
                \texttt{That sounds like a good starting point, but I'm not sure if I'm willing to give up the pepper. It's worth something to me, and I'm not sure if the hammer and 2 ropes are enough to make up for it. } \\
\\ 
\texttt{Can you tell me what you think the pepper is worth to you, and if you'd be willing to trade for it without giving up the hammer?} \\
            }
        }
    }
    & & \\ \\

    \theutterance \stepcounter{utterance}  
    & & \multicolumn{4}{p{0.6\linewidth}}{
        \cellcolor[rgb]{0.9,0.9,0.9}{
            \makecell[{{p{\linewidth}}}]{
                \texttt{\tiny{[P1$\langle$GM]}}
                \texttt{That sounds like a good starting point, but I'm not sure if I'm willing to give up the pepper. It's worth something to me, and I'm not sure if the hammer and 2 ropes are enough to make up for it. } \\
\\ 
\texttt{Can you tell me what you think the pepper is worth to you, and if you'd be willing to trade for it without giving up the hammer?} \\
            }
        }
    }
    & \\ \\

    \theutterance \stepcounter{utterance}  
    & \multicolumn{4}{p{0.6\linewidth}}{
        \cellcolor[rgb]{0.9,0.9,0.9}{
            \makecell[{{p{\linewidth}}}]{
                \texttt{\tiny{[P1$\rangle$GM]}}
                \texttt{The pepper is worth 1 point to me, but I'm trying to maximize my score. I'm willing to trade for it, but I'd like to get a better deal. } \\
\\ 
\texttt{How about this: I'll give you the hammer and 2 ropes in exchange for the pepper and 1 rope, but I'll also throw in 1 additional rope?} \\
            }
        }
    }
    & & \\ \\

    \theutterance \stepcounter{utterance}  
    & & & \multicolumn{2}{p{0.3\linewidth}}{
        \cellcolor[rgb]{0.9,0.9,0.9}{
            \makecell[{{p{\linewidth}}}]{
                \texttt{\tiny{[GM$|$GM]}}
                \texttt{continue} \\
            }
        }
    }
    & & \\ \\

    \theutterance \stepcounter{utterance}  
    & & & \multicolumn{2}{p{0.3\linewidth}}{
        \cellcolor[rgb]{0.9,0.9,0.9}{
            \makecell[{{p{\linewidth}}}]{
                \texttt{\tiny{[GM$|$GM]}}
                \texttt{The pepper is worth 1 point to me, but I'm trying to maximize my score. I'm willing to trade for it, but I'd like to get a better deal. } \\
\\ 
\texttt{How about this: I'll give you the hammer and 2 ropes in exchange for the pepper and 1 rope, but I'll also throw in 1 additional rope?} \\
            }
        }
    }
    & & \\ \\

    \theutterance \stepcounter{utterance}  
    & & \multicolumn{4}{p{0.6\linewidth}}{
        \cellcolor[rgb]{0.9,0.9,0.9}{
            \makecell[{{p{\linewidth}}}]{
                \texttt{\tiny{[GM$\rangle$P2]}}
                \texttt{The pepper is worth 1 point to me, but I'm trying to maximize my score. I'm willing to trade for it, but I'd like to get a better deal. } \\
\\ 
\texttt{How about this: I'll give you the hammer and 2 ropes in exchange for the pepper and 1 rope, but I'll also throw in 1 additional rope?} \\
            }
        }
    }
    & \\ \\

    \theutterance \stepcounter{utterance}  
    & & & \multicolumn{4}{p{0.6\linewidth}}{
        \cellcolor[rgb]{0.9,0.9,0.9}{
            \makecell[{{p{\linewidth}}}]{
                \texttt{\tiny{[GM$\langle$P2]}}
                \texttt{That's a good offer, but I'm not sure if I'm willing to give up the pepper for just 1 additional rope. I think the pepper is worth more to me than that. } \\
\\ 
\texttt{Can you tell me what you think the hammer is worth to you, and if you'd be willing to trade for it without giving up any of the ropes?} \\
            }
        }
    }
     \\ \\

    \theutterance \stepcounter{utterance}  
    & & & \multicolumn{2}{p{0.3\linewidth}}{
        \cellcolor[rgb]{0.9,0.9,0.9}{
            \makecell[{{p{\linewidth}}}]{
                \texttt{\tiny{[GM$|$GM]}}
                \texttt{continue} \\
            }
        }
    }
    & & \\ \\

    \theutterance \stepcounter{utterance}  
    & & & \multicolumn{2}{p{0.3\linewidth}}{
        \cellcolor[rgb]{0.9,0.9,0.9}{
            \makecell[{{p{\linewidth}}}]{
                \texttt{\tiny{[GM$|$GM]}}
                \texttt{That's a good offer, but I'm not sure if I'm willing to give up the pepper for just 1 additional rope. I think the pepper is worth more to me than that. } \\
\\ 
\texttt{Can you tell me what you think the hammer is worth to you, and if you'd be willing to trade for it without giving up any of the ropes?} \\
            }
        }
    }
    & & \\ \\

    \theutterance \stepcounter{utterance}  
    & & \multicolumn{4}{p{0.6\linewidth}}{
        \cellcolor[rgb]{0.9,0.9,0.9}{
            \makecell[{{p{\linewidth}}}]{
                \texttt{\tiny{[P1$\langle$GM]}}
                \texttt{That's a good offer, but I'm not sure if I'm willing to give up the pepper for just 1 additional rope. I think the pepper is worth more to me than that. } \\
\\ 
\texttt{Can you tell me what you think the hammer is worth to you, and if you'd be willing to trade for it without giving up any of the ropes?} \\
            }
        }
    }
    & \\ \\

    \theutterance \stepcounter{utterance}  
    & \multicolumn{4}{p{0.6\linewidth}}{
        \cellcolor[rgb]{0.9,0.9,0.9}{
            \makecell[{{p{\linewidth}}}]{
                \texttt{\tiny{[P1$\rangle$GM]}}
                \texttt{The hammer is worth 3 points to me, and it's a very valuable item for me. I'm willing to trade for it, but I'd like to get a good deal. } \\
\\ 
\texttt{How about this: I'll give you the hammer in exchange for 2 ropes?} \\
            }
        }
    }
    & & \\ \\

    \theutterance \stepcounter{utterance}  
    & & & \multicolumn{2}{p{0.3\linewidth}}{
        \cellcolor[rgb]{0.9,0.9,0.9}{
            \makecell[{{p{\linewidth}}}]{
                \texttt{\tiny{[GM$|$GM]}}
                \texttt{continue} \\
            }
        }
    }
    & & \\ \\

    \theutterance \stepcounter{utterance}  
    & & & \multicolumn{2}{p{0.3\linewidth}}{
        \cellcolor[rgb]{0.9,0.9,0.9}{
            \makecell[{{p{\linewidth}}}]{
                \texttt{\tiny{[GM$|$GM]}}
                \texttt{The hammer is worth 3 points to me, and it's a very valuable item for me. I'm willing to trade for it, but I'd like to get a good deal. } \\
\\ 
\texttt{How about this: I'll give you the hammer in exchange for 2 ropes?} \\
            }
        }
    }
    & & \\ \\

    \theutterance \stepcounter{utterance}  
    & & \multicolumn{4}{p{0.6\linewidth}}{
        \cellcolor[rgb]{0.9,0.9,0.9}{
            \makecell[{{p{\linewidth}}}]{
                \texttt{\tiny{[GM$\rangle$P2]}}
                \texttt{The hammer is worth 3 points to me, and it's a very valuable item for me. I'm willing to trade for it, but I'd like to get a good deal. } \\
\\ 
\texttt{How about this: I'll give you the hammer in exchange for 2 ropes?} \\
            }
        }
    }
    & \\ \\

    \theutterance \stepcounter{utterance}  
    & & & \multicolumn{4}{p{0.6\linewidth}}{
        \cellcolor[rgb]{0.9,0.9,0.9}{
            \makecell[{{p{\linewidth}}}]{
                \texttt{\tiny{[GM$\langle$P2]}}
                \texttt{That's a good offer, but I'm not sure if I'm willing to give up 2 ropes for the hammer. I think the ropes are worth more to me than that. } \\
\\ 
\texttt{I'll tell you what, I'll give you the hammer in exchange for 1 rope and the pepper. That way, I get the pepper, which is worth 1 point to me, and you get the hammer, which is worth 3 points to you.} \\
            }
        }
    }
     \\ \\

    \theutterance \stepcounter{utterance}  
    & & & \multicolumn{2}{p{0.3\linewidth}}{
        \cellcolor[rgb]{0.9,0.9,0.9}{
            \makecell[{{p{\linewidth}}}]{
                \texttt{\tiny{[GM$|$GM]}}
                \texttt{continue} \\
            }
        }
    }
    & & \\ \\

    \theutterance \stepcounter{utterance}  
    & & & \multicolumn{2}{p{0.3\linewidth}}{
        \cellcolor[rgb]{0.9,0.9,0.9}{
            \makecell[{{p{\linewidth}}}]{
                \texttt{\tiny{[GM$|$GM]}}
                \texttt{That's a good offer, but I'm not sure if I'm willing to give up 2 ropes for the hammer. I think the ropes are worth more to me than that. } \\
\\ 
\texttt{I'll tell you what, I'll give you the hammer in exchange for 1 rope and the pepper. That way, I get the pepper, which is worth 1 point to me, and you get the hammer, which is worth 3 points to you.} \\
            }
        }
    }
    & & \\ \\

    \theutterance \stepcounter{utterance}  
    & & \multicolumn{4}{p{0.6\linewidth}}{
        \cellcolor[rgb]{0.9,0.9,0.9}{
            \makecell[{{p{\linewidth}}}]{
                \texttt{\tiny{[P1$\langle$GM]}}
                \texttt{That's a good offer, but I'm not sure if I'm willing to give up 2 ropes for the hammer. I think the ropes are worth more to me than that. } \\
\\ 
\texttt{I'll tell you what, I'll give you the hammer in exchange for 1 rope and the pepper. That way, I get the pepper, which is worth 1 point to me, and you get the hammer, which is worth 3 points to you.} \\
            }
        }
    }
    & \\ \\

    \theutterance \stepcounter{utterance}  
    & \multicolumn{4}{p{0.6\linewidth}}{
        \cellcolor[rgb]{0.9,0.9,0.9}{
            \makecell[{{p{\linewidth}}}]{
                \texttt{\tiny{[P1$\rangle$GM]}}
                \texttt{That's a fair deal. I think I can agree to that. } \\
\\ 
\texttt{{[}Proposal: 1 rope, 1 pepper, hammer{]}} \\
            }
        }
    }
    & & \\ \\

    \theutterance \stepcounter{utterance}  
    & & & \multicolumn{2}{p{0.3\linewidth}}{
        \cellcolor[rgb]{0.9,0.9,0.9}{
            \makecell[{{p{\linewidth}}}]{
                \texttt{\tiny{[GM$|$GM]}}
                \texttt{proposal does not include number/name pairs} \\
            }
        }
    }
    & & \\ \\

    \theutterance \stepcounter{utterance}  
    & \multicolumn{4}{p{0.6\linewidth}}{
        \cellcolor[rgb]{0.95,0.95,0.95}{
            \makecell[{{p{\linewidth}}}]{
                \texttt{\tiny{[P2$\rightarrow$P1]}}
                \texttt{Post{-}episode Reflexion (not sent to GM):} \\
\texttt{New plan: } \\
\\ 
\texttt{In my last attempt, I was able to make a trade that seemed beneficial, but ultimately, I failed to achieve the target state of the game. Upon reflection, I realize that I didn't consider the other player's secret value function and my own value function in a more strategic way. } \\
\\ 
\texttt{In the beginning, I tried to gauge the other player's interest in the hammer, which is a valuable item for me. However, I didn't use this information effectively to make a proposal that would maximize my score. } \\
\\ 
\texttt{Moreover, I made a mistake by not making a secret proposal at the end of the game. I should have proposed a deal that would have maximized my score, rather than agreeing to the other player's proposal. } \\
\\ 
\texttt{Next time, I will start by making a proposal that takes into account both my own value function and the other player's interest in the items. I will also try to identify the items that are most valuable to both of us and make a trade that benefits both parties. } \\
\\ 
\texttt{To avoid making a secret proposal too late, I will make sure to keep track of the items that are left and the value of each item. This way, I can make a proposal that maximizes my score at the end of the game. } \\
\\ 
\texttt{For example, I could propose giving the other player the items that are most valuable to them, in exchange for the items that are most valuable to me. This way, I would get the items that I need to maximize my score, and the other player would get the items that they need to maximize their score. } \\
\\ 
\texttt{By considering the other player's secret value function and making a proposal that benefits both parties, I hope to achieve a higher score and win the game.} \\
            }
        }
    }
    & & \\ \\

    \theutterance \stepcounter{utterance}  
    & & & \multicolumn{4}{p{0.6\linewidth}}{
        \cellcolor[rgb]{0.95,0.95,0.95}{
            \makecell[{{p{\linewidth}}}]{
                \texttt{\tiny{[P1$\rightarrow$P2]}}
                \texttt{Post{-}episode Reflexion (not sent to GM):} \\
\texttt{New plan: } \\
\\ 
\texttt{In my last attempt, I failed to maximize my score because I didn't consider the value function of the other player and made a secret proposal that was not complementary to theirs. I should have continued to negotiate and find a mutually beneficial agreement.} \\
\\ 
\texttt{To improve, I will continue to ask open{-}ended questions to gauge the other player's interest in each item and understand their value function. I will also pay closer attention to the other player's responses and adjust my strategy accordingly.} \\
\\ 
\texttt{In particular, I should have noticed that the other player was willing to trade for the hammer and ropes, but was not willing to give up the pepper. I should have taken this into account when making my secret proposal and included the pepper in the trade.} \\
\\ 
\texttt{Next time, I will make sure to continue negotiating and find a mutually beneficial agreement that takes into account the value function of both players. I will also make sure to follow the rules and format of the game, avoiding any secret proposals that are not complementary to the other player's proposal.} \\
            }
        }
    }
     \\ \\

\end{supertabular}
}

\end{document}
